\section{Results}
\label{sec:results}

\subsection{Integrity and hypothesis outcomes}

All 800 evaluation queries were retained.  Independent BFS and deterministic
search replay passed for every query; all four methods returned the BFS cost,
and \EagerFull{}, \LazyFull{}, and \Staged{} had matching full-landmark
expansion digests in all 800 cases.  No post-hoc exclusion was made.  These
checks passed before timing was analyzed.

\begin{table*}[t]
\centering
\small
\caption{Outcomes under the prespecified analysis.  H3 is descriptive;
its bootstrap fraction and correlation are not significance or causal tests.}
\label{tab:hypotheses}
\resizebox{\textwidth}{!}{\begin{tabular}{llll}
\toprule
Hypothesis & Status & Estimate & Interval \\
\midrule
H1 & passed & 800/800 queries passed both gates & not applicable (exact validation) \\
H2 & supported & pivot=0.0309758772; read=0.0301724138 & descriptive IQRs in summary/map tables \\
H3 & descriptive & ratio=1.05430023; Spearman rho=-0.640610853 & hierarchical bootstrap 95\% CI [0.989043253, 1.10158524] \\
\bottomrule
\end{tabular}
}
\end{table*}

Table~\ref{tab:hypotheses} reports the prespecified outcomes.  H2 was supported:
relative to \LazyFull{}, \Staged{} saved a median 3.0976\% of pivot
evaluations and 3.0172\% of table reads.  The corresponding query-level
medians were 28 fewer pivots/reads and one avoided suffix call, while the
median additional OPEN work was 116 combined queue operations (insertions and
pops).  Pivot-evaluation savings were nonnegative (420 positive, 380 zero) and
skewed: the median was 3.10\%, the mean 13.44\%, and the IQR 0--27.34\%.  This is a
heuristic-work improvement, not a wall-clock improvement.

\subsection{Search time and deterministic work}

The primary staged/lazy time ratio was 1.054300.  Its map-then-query
hierarchical bootstrap interval was $[0.989043,1.101585]$; 5.3\% of bootstrap
replicates had a ratio below one.  Since the interval includes one and only
eight top-level clusters are available, we make no categorical dominance or
significance claim.

\begin{table}[t]
\centering
\scriptsize
\caption{Overall per-query summaries on 800 evaluation queries.  Time is the
mean of each query's eight-run median, in ms; counters are query means.}
\label{tab:methods}
\begin{tabular}{lrrrrr}
\toprule
Method & Time & Expanded & Pops & Requeues & Reads \\
\midrule
Manhattan & 15.535 & 1197.86 & 1356.08 & 0.00 & 0.00 \\
Eager full & 8.149 & 160.69 & 162.09 & 0.00 & 9977.68 \\
Lazy full & 8.282 & 160.69 & 422.15 & 260.06 & 8385.88 \\
Staged & 8.056 & 160.69 & 599.86 & 437.78 & 6080.29 \\
\bottomrule
\end{tabular}
\end{table}

Table~\ref{tab:methods} exposes the tradeoff.  All full-landmark schedules
averaged 160.69 expansions and 517.30 generated successors per query,
compared with 1197.86 and 4043.00 for Manhattan.  Eager evaluation minimized
OPEN operations but averaged 9,945.68 pivot evaluations.  Ordinary lazy
evaluation reduced this to 8,353.88 at the cost of 260.06 requeues.  Staging
reduced it further to 6,048.29 but averaged 437.78 requeues and 599.86 pops.
No method reopened a state.  Table~\ref{tab:methods} averages query medians,
whereas the primary statistic is a median of paired time ratios.  These are
different estimands, especially under strong map heterogeneity.

\subsection{Map and family crossover}

\begin{figure}[t]
  \centering
  \includegraphics[width=\columnwidth]{per_map_time_ratios.pdf}
  \caption{Per-map distributions of paired staged/lazy search-time ratios.
  The horizontal reference at one denotes equal time.}
  \label{fig:map-ratios}
\end{figure}

\begin{table}[t]
\centering
\scriptsize
\caption{Per-map medians.  Ratio $<1$ favors staged evaluation; read saving is
relative to ordinary lazy evaluation.  V and H denote validation- and
holdout-source maps.}
\label{tab:maps}
\begin{tabular}{llrr}
\toprule
Family & Source/map & Time ratio & Read saving \\
\midrule
Maze & V: 128-2 & 0.935 & 31.17\% \\
Maze & H: 128-10 & 0.934 & 30.34\% \\
Random & V: 32-20 & 1.069 & 2.50\% \\
Random & H: 32-10 & 1.110 & 0.00\% \\
Room & V: 64-16 & 0.990 & 18.47\% \\
Room & H: 32-4 & 1.039 & 11.46\% \\
Warehouse & V: 10-20-1 & 1.137 & 0.00\% \\
Warehouse & H: 10-20-2 & 1.157 & 0.00\% \\
\bottomrule
\end{tabular}
\end{table}

Figure~\ref{fig:map-ratios} and Table~\ref{tab:maps} show repeatable family
regimes.  Both maze maps saved about 30--31\% of reads and favored staging by
about 6.5\%.  Both warehouse maps had zero median read saving and favored
ordinary lazy evaluation by 13.7--15.7\%.  Random maps favored lazy, while the
two room maps lay on opposite sides of one.  The same direction on the two
maze and two warehouse maps is suggestive replication within those families,
not a population-level guarantee.

\begin{figure*}[t]
  \centering
  \includegraphics[width=0.92\textwidth]{family_mechanism_decomposition.pdf}
  \caption{Family-level decomposition of saved landmark work and added OPEN
  work.  Families are equal-weighted through their constituent queries.}
  \label{fig:families}
\end{figure*}

\begin{figure*}[t]
  \centering
  \includegraphics[width=0.92\textwidth]{saved_work_vs_time.pdf}
  \caption{Query-level exact work changes versus paired time ratio.  The
  association is descriptive and queries within a map are dependent.}
  \label{fig:saved-work}
\end{figure*}

Across all 800 queries, read-saving fraction and log time ratio had Spearman
$\rho=-0.640611$; extra OPEN work and log ratio had $\rho=-0.240994$.
Figure~\ref{fig:saved-work} makes the crossover visible but does not isolate a
causal coefficient: difficulty, topology, and absolute search size affect both
the counters and time.

\subsection{Preprocessing amortization}

\begin{figure}[t]
  \centering
  \includegraphics[width=\columnwidth]{preprocessing_amortization.pdf}
  \caption{Equal-map-weighted amortized time as the one-time 32-landmark table
  construction is spread over $Q$ queries.}
  \label{fig:amortization}
\end{figure}

\begin{table}[t]
\centering
\scriptsize
\caption{Mean amortized ms/query, equal-weighting eight maps.}
\label{tab:amortization}
\begin{tabular}{rrrrr}
\toprule
$Q$ & Manhattan & Eager & Lazy & Staged \\
\midrule
1 & 15.535 & 1881.862 & 1881.995 & 1881.769 \\
10 & 15.535 & 195.520 & 195.653 & 195.427 \\
100 & 15.535 & 26.886 & 27.019 & 26.793 \\
1000 & 15.535 & 10.023 & 10.155 & 9.929 \\
\bottomrule
\end{tabular}
\end{table}

Preprocessing dominates when a table serves few queries
(Figure~\ref{fig:amortization}, Table~\ref{tab:amortization}).  At $Q=100$ all
landmark modes remain slower than Manhattan after charging construction;
at $Q=1000$ they are faster under this equal-map average.  Because all three
landmark schedules share exactly the same table, preprocessing barely changes
their relative order; it changes the practical comparison with Manhattan.
